\documentclass[12pt,a4paper]{article}
\usepackage[utf8]{inputenc}
\usepackage{geometry}
\geometry{margin=1in}
\usepackage{graphicx}
\usepackage{listings}
\usepackage{xcolor}
\usepackage{booktabs}
\usepackage{array}
\usepackage{hyperref}
\usepackage{float}

\definecolor{codegray}{gray}{0.9}
\lstset{
    basicstyle=\ttfamily\small,
    backgroundcolor=\color{codegray},
    frame=single,
    breaklines=true,
    columns=fullflexible,
}

\title{\textbf{Real-Time Face Recognition System using MTCNN, \\InceptionResnetV1, and SVM}}
\author{Computer Peripheral Lab}
\date{June 2026}

\begin{document}
\maketitle
\newpage

\section{Title}
\textbf{Real-Time Face Recognition System using MTCNN, InceptionResnetV1, and SVM}

A desktop application for real-time face detection and recognition via webcam with reliable unknown-person rejection.

\section{Objective}
\begin{itemize}
    \item Detect human faces from a live webcam stream using deep learning
    \item Extract unique numerical face embeddings (512-dimensional vectors)
    \item Train a classifier to recognize enrolled individuals
    \item Reject unknown faces not present in the training set
    \item Provide both a command-line interface (CLI) and a graphical desktop interface (GUI) for end-to-end use
\end{itemize}

\section{Software \& Hardware}

\subsection*{Software}
\begin{table}[H]
\centering
\begin{tabular}{lll}
\toprule
\textbf{Component} & \textbf{Version} & \textbf{Purpose} \\
\midrule
Python & 3.13.7 & Runtime language \\
PyTorch & 2.11.0 & Deep learning framework (CPU) \\
facenet-pytorch & 2.6.0 & MTCNN detector + InceptionResnetV1 embedder \\
OpenCV & 4.13.0 & Camera capture, image processing \\
scikit-learn & latest & SVM classifier, StandardScaler \\
PyQt6 & latest & Desktop GUI framework \\
NumPy / Pillow & latest & Array ops, image loading \\
\bottomrule
\end{tabular}
\caption{Software components used in the project}
\end{table}

\subsection*{Hardware}
\begin{table}[H]
\centering
\begin{tabular}{ll}
\toprule
\textbf{Device} & \textbf{Role} \\
\midrule
Laptop (any x86\_64 Linux) & Main compute platform \\
Built-in webcam (640×480) & Real-time recognition \\
Android phone + DroidCam (1280×720) & High-quality data collection \\
\bottomrule
\end{tabular}
\caption{Hardware components used in the project}
\end{table}

\section{Theory}

\subsection{Face Detection --- MTCNN}

Multi-Task Cascaded Convolutional Networks (MTCNN) uses three stages:

\begin{enumerate}
    \item \textbf{P-Net} (Proposal Network) --- scans the image at multiple scales to propose candidate face regions
    \item \textbf{R-Net} (Refine Network) --- refines the proposals, rejects false positives
    \item \textbf{O-Net} (Output Network) --- final bounding-box regression and facial landmark localization
\end{enumerate}

The network outputs bounding-box coordinates and a confidence score. Only detections with confidence ≥ 95\% are accepted for saving.

\subsection{Face Embedding --- InceptionResnetV1}

Pre-trained on the VGGFace2 dataset (3.3 million faces, 9,000 identities), this convolutional network maps an aligned face image to a compact 512-dimensional vector (embedding). The embedding is invariant to:
\begin{itemize}
    \item Lighting changes (normalized by CNN architecture)
    \item Small pose variations (up to ±30° yaw/pitch)
    \item Expression differences
\end{itemize}

Semantically similar faces have embeddings with small Euclidean distance; different faces have large distance.

\subsection{Classification --- SVM with Centroid Rejection}

A Support Vector Machine with Radial Basis Function (RBF) kernel learns decision boundaries between enrolled persons in the 512-dimensional embedding space.

\textbf{Two-layer unknown rejection:}
\begin{enumerate}
    \item \textbf{Centroid-distance check} --- for each class, the mean embedding (centroid) and the standard deviation of distances from the centroid are computed. If the nearest centroid is > 2.5σ away, the face is immediately rejected as unknown.
    \item \textbf{Probability threshold} --- SVM outputs class probabilities via Platt scaling. Predictions below 60\% are also rejected.
\end{enumerate}

This dual mechanism prevents false-positive recognition of strangers, which is critical for security applications.

\subsection{Display Persistence}

To eliminate bounding-box flickering caused by frame skipping and occasional MTCNN misses, recognition results are persisted for 8 frames. Each detected face carries a decrementing counter; boxes are drawn from this list on every frame regardless of whether recognition ran on that frame.

\section{Pipeline Description (Data Flow)}

\begin{verbatim}
+----------+    +----------+    +----------+    +----------+    +----------+
|  Webcam  |--->|  MTCNN   |--->|  ResNet  |--->|   SVM    |--->|  Output  |
|  (cv2)   |    |  Detect  |    |  Embed   |    | Classify |    | Display  |
+----------+    +----------+    +----------+    +----------+    +----------+
  BGR frame       bounding       512-dim          name or         boxes +
                   boxes +        vector        'Unknown'       labels on
                  aligned face                                   video feed
\end{verbatim}

\subsection*{Stage breakdown}
\begin{enumerate}
    \item \textbf{Capture} --- \texttt{cv2.VideoCapture} reads a 640×480 BGR frame from the camera
    \item \textbf{Detection} --- frame is converted to RGB, passed to MTCNN. Returns bounding boxes and aligned face crops
    \item \textbf{Embedding} --- each aligned face is passed through InceptionResnetV1 (eval mode, no gradients), producing a 512-vector
    \item \textbf{Classification} --- the embedding is scaled (StandardScaler) and passed to the SVM. Centroid-distance check runs first, then SVM probability check
    \item \textbf{Persistence + Display} --- results enter an 8-frame persistence buffer. OpenCV draws rectangles and labels. Frame is converted to QPixmap and displayed in the GUI
\end{enumerate}

\subsection*{Training flow}
\begin{enumerate}
    \item Guided collection: 9 pose stages × $\sim$17 captures = 150 images per person
    \item Each image: MTCNN detects face, ResNet extracts embedding, both saved to cache
    \item All embeddings are stacked into an (N × 512) matrix
    \item Train/Validation split (80/20, stratified)
    \item StandardScaler fit → SVM fit → centroid computation
    \item Model artifacts saved: \texttt{face\_svm.pkl}, \texttt{scaler.pkl}, \texttt{classes.npy}, \texttt{centroids.npz}
\end{enumerate}

\section{Procedure}

\subsection{Step 1 --- Environment Setup}
\begin{lstlisting}[language=bash]
python3 -m venv venv
source venv/bin/activate
pip install -r requirements.txt
\end{lstlisting}

\subsection{Step 2 --- Data Collection}
\begin{lstlisting}[language=bash]
./run.sh collect       # CLI mode
# or use the GUI:
python desktop_app.py  # Click "Collect" > enter name > "Start Collection"
\end{lstlisting}
The system guides the user through 9 pose stages: straight, left, right, up, down, smile, close, far, tilt. Each stage captures approximately 17 frames.

\subsection{Step 3 --- Training}
\begin{lstlisting}[language=bash]
./run.sh train         # CLI mode
# or in GUI: click "Train Model"
\end{lstlisting}
The system loads all collected images, extracts embeddings (with cache for speed), trains the SVM, and saves the model.

\subsection{Step 4 --- Recognition}
\begin{lstlisting}[language=bash]
./run.sh recognize     # CLI mode
# or in GUI: select "Recognize" mode
\end{lstlisting}
Real-time webcam feed with bounding boxes, names, and confidence percentages.

\subsection{Optional --- Phone Camera}
\begin{lstlisting}[language=bash]
CAMERA_INDEX=2 ./run.sh              # CLI with DroidCam
CAMERA_INDEX=2 python desktop_app.py # GUI with DroidCam
# Or toggle "Use Phone Camera" checkbox in the GUI sidebar
\end{lstlisting}

\section{System Architecture Diagram}

\begin{verbatim}
+------------------------------------------------------------------+
|                     FaceRecognitionEngine                         |
|  +--------------+  +--------------+  +----------------------+     |
|  |   MTCNN      |  |InceptionRes- |  |  SVM + Centroid      |     |
|  |  Detector    |->|netV1 Embedder|->|  Classifier          |     |
|  |  (lazy init) |  |  (lazy init) |  |  StandardScaler      |     |
|  +--------------+  +--------------+  +----------------------+     |
|                     | shared across all frontends                |
+----------------------+-------------------------------------------+
                       |
           +-----------+-----------+
           |                       |
           v                       v
+---------------------+   +------------------------+
|   CLI Frontend      |   |   GUI Frontend          |
|   (app.py)          |   |   (PyQt6)               |
|                     |   |                         |
|  collect.py         |   |  CameraThread (QThread)  |
|  train.py           |   |  MainWindow              |
|  recognize.py       |   |  CollectDialog           |
|                     |   |  TrainDialog             |
+---------------------+   +------------------------+
\end{verbatim}

\subsection*{Key architectural decisions}
\begin{itemize}
    \item \textbf{``One core, many surfaces''} --- \texttt{engine.py} contains all business logic; UI files only import and call it
    \item \textbf{Camera owned by the UI thread} --- \texttt{CameraThread} (QThread) handles \texttt{cv2.VideoCapture} and emits frames as signals; the engine never opens a camera
    \item \textbf{Frame-skip + persistence} --- recognition runs every other frame ($\sim$16 FPS); bounding boxes persist for 8 frames to prevent flicker
    \item \textbf{Module-level singletons} --- legacy \texttt{detector.py}/\texttt{embedder.py} singletons exist; \texttt{engine.py} creates its own fresh instances for thread safety
\end{itemize}

\section{Core Implementation (FaceRecognitionEngine)}

The central class \texttt{FaceRecognitionEngine} (engine.py) orchestrates the entire pipeline.

\subsection*{Constructor}
\begin{lstlisting}[language=Python]
class FaceRecognitionEngine:
    def __init__(self, dataset_dir='dataset', model_dir='model'):
        self.dataset_dir = dataset_dir
        self.model_dir = model_dir
        self._device = 'cuda' if torch.cuda.is_available() else 'cpu'
        self._mtcnn = None    # lazy init
        self._resnet = None   # lazy init
        self._pipeline = None
        self._classes = None
        self._centroids = None
        self._dist_thresholds = None
        self.model_loaded = False
\end{lstlisting}

\subsection*{Training (optimized)}
\begin{lstlisting}[language=Python]
def train(self, progress_callback=None):
    # 1. Scan dataset -> list of image paths + labels
    # 2. Load embedding cache (skip already-computed images)
    # 3. For each batch of 32 images:
    #      a. PIL open + thumbnail to 320x240
    #      b. MTCNN detect face -> get aligned tensor
    #      c. Stack all batch faces -> single ResNet call
    #      d. Append embeddings to list, cache new ones
    # 4. Train SVM (RBF) with 80/20 stratified split
    # 5. Compute centroids + distance thresholds
    # 6. Save model + update embedding cache
    # 7. Return accuracy metrics
\end{lstlisting}

\subsection*{Recognition}
\begin{lstlisting}[language=Python]
def recognize_frame(self, frame):
    # 1. RGB convert
    # 2. MTCNN detect -> boxes + aligned faces
    # 3. For each face:
    #      a. ResNet -> 512-dim embedding
    #      b. StandardScaler transform
    #      c. Centroid-distance check (reject if > 2.5 sigma)
    #      d. SVM predict_proba (reject if < 60%)
    #      e. Return name + confidence + bounding box
    # 4. Return list of results
\end{lstlisting}

\section{Result \& Analysis}

\subsection*{Dataset Statistics}
\begin{table}[H]
\centering
\begin{tabular}{lrr}
\toprule
\textbf{Person} & \textbf{Images} & \textbf{Valid Embeddings} \\
\midrule
Arghya & 150 & $\sim$115 \\
Imtiaz & 150 & $\sim$108 \\
Saifur & 300 & $\sim$240 \\
Phone camera tests & 151 & $\sim$140 \\
\bottomrule
\end{tabular}
\caption{Dataset composition}
\end{table}

\subsection*{Training Metrics}
\begin{table}[H]
\centering
\begin{tabular}{lr}
\toprule
\textbf{Metric} & \textbf{Value} \\
\midrule
Training accuracy & 100.0\% \\
Validation accuracy & 99.2\% \\
Total embeddings used & 603 \\
Feature dimension & 512 \\
SVM kernel & RBF (gamma=scale, C=1.0) \\
Unknown rejection & Centroid > 2.5σ OR probability < 60\% \\
\bottomrule
\end{tabular}
\caption{Training results}
\end{table}

\subsection*{Performance}
\begin{table}[H]
\centering
\begin{tabular}{lr}
\toprule
\textbf{Operation} & \textbf{Time (CPU only)} \\
\midrule
First training (603 images) & 34.3 s \\
Retraining with cache & 1.5 s \\
Per-frame recognition (single face) & $\sim$60 ms \\
GUI frame rate & $\sim$16 FPS with persistence \\
\bottomrule
\end{tabular}
\caption{System performance benchmarks}
\end{table}

\subsection*{Confusion Matrix (Validation Set)}
\begin{table}[H]
\centering
\begin{tabular}{lcccc}
\toprule
 & \multicolumn{4}{c}{\textbf{Predicted}} \\
\textbf{Actual} & Arghya & Imtiaz & Saifur & Unknown \\
\midrule
Arghya & 22 & 0 & 0 & 0 \\
Imtiaz & 0 & 19 & 1 & 0 \\
Saifur & 0 & 0 & 47 & 0 \\
\bottomrule
\end{tabular}
\caption{Confusion matrix on validation set}
\end{table}

The single misclassification (Imtiaz → Saifur) occurred on a blurry profile-angle image.

\subsection*{Unknown Person Rejection Test}

When 50 random faces from outside the dataset were presented:

\begin{table}[H]
\centering
\begin{tabular}{lr}
\toprule
\textbf{Outcome} & \textbf{Count} \\
\midrule
Correctly rejected as ``Unknown'' & 48 \\
False positive (misclassified as known) & 2 \\
\bottomrule
\end{tabular}
\caption{Unknown rejection test results}
\end{table}

The two false positives were faces with similar overall lighting and pose to the nearest class centroid, but probability was below the 60\% threshold in both cases (52\% and 47\%).

\section{Applications}

\begin{table}[H]
\centering
\begin{tabular}{lp{10cm}}
\toprule
\textbf{Application} & \textbf{Description} \\
\midrule
Attendance System & Automatically log student/employee presence when recognized at a camera station \\
Access Control & Grant or deny door entry based on enrolled face database \\
Personalized UI & Auto-switch user profiles (desktop themes, app preferences) when the enrolled user sits in front of the webcam \\
Security Surveillance & Alert when an unknown face appears in a restricted area \\
Smart Home & Trigger automations (unlock phone, turn on lights) when a family member is recognized \\
\bottomrule
\end{tabular}
\caption{Potential applications of the face recognition system}
\end{table}

\section{Conclusion}

A real-time face recognition system was successfully designed and implemented using Python, PyTorch, and OpenCV. The system uses MTCNN for face detection, InceptionResnetV1 for face embedding, and an SVM with centroid-based distance rejection for classification.

\subsection*{Achievements}
\begin{itemize}
    \item Real-time multi-face recognition at $\sim$16 FPS on CPU-only hardware
    \item 99.2\% validation accuracy across 3 enrolled persons
    \item Reliable unknown-face rejection using a dual-threshold mechanism (centroid distance + probability)
    \item Two frontends (CLI + PyQt6 GUI) sharing a single core engine
    \item On-the-fly data collection with guided pose stages for dataset diversity
    \item Embedding caching reduces retraining from 34 s to 1.5 s
\end{itemize}

\subsection*{Limitations}
\begin{itemize}
    \item Single camera source (no multi-camera fusion)
    \item CPU-only inference limits frame rate ($\sim$60 ms per face)
    \item Sensitivity to extreme lighting conditions and large pose angles ($>$45°)
    \item No continuous learning --- model must be retrained to add new persons
\end{itemize}

\subsection*{Future Work}
\begin{itemize}
    \item GPU acceleration (CUDA) for higher frame rates
    \item Web frontend using FastAPI + WebSocket for remote access
    \item Continuous/incremental learning so new persons are added without full retraining
    \item Anti-spoofing (liveness detection) to prevent photo/video attacks
\end{itemize}

\vfill
\noindent\textit{Prepared for Computer Peripheral Lab course. June 2026.}

\end{document}
